\paragraph*{In short}

You are now aware that a Croqtile program consists of three parts, with the \emph{tileflow program} being the core. This part is transpiled into target source code during the compilation process. The call chain typically flows from the host to the tileflow program, and then from the tileflow program to the device code.

You may have noticed that Croqtile not only abstracts DMA operations into higher-level semantics but also combines iteration and tiling for ease of use. This makes Croqtile code concise and expressive. In the following chapters, you will delve deeper into Croqtile's syntax and semantics to explore its full potential.
